\section*{Chapter \ref{ch:hf:on-chain-trading} --- On-Chain Trading}

\begin{solution}{exo:hf:on-chain-trading:1}
The pool's price, \$3\,100, is above \$3\,000 by more than the fee: sell
\[
\Delta x^\ast=\frac{\sqrt{0.997\times1\,000\times3\,100\,000/3\,000}-1\,000}{0.997}=15.05
\]
tokens, for a profit of \$677 before gas.
\end{solution}

\begin{solution}{exo:hf:on-chain-trading:2}
$0.04^2/8\times20\,000\,000=\$4\,000$ a day, \$28\,000 a week.
\end{solution}

\begin{solution}{exo:hf:on-chain-trading:3}
$1\,000\,000\times0.05-20=\$49\,980$ before the \omterm{def:hf:on-chain-trading:bid}{bundle bid}.
\end{solution}

\begin{solution}{exo:hf:on-chain-trading:4}
Each searcher's value differs from the others' only by its costs; with more competitors the second-best value approaches the best, and in a
first-price auction the winner must bid close to it: what it keeps is only its cost advantage.
\end{solution}

\begin{solution}{exo:hf:on-chain-trading:5}
Loss-versus-rebalancing is set by how far the market price moves between arbitrages, which depends on volatility and the pool's size; the fee only
widens the band inside which arbitrage does not pay, which changes the loss slightly.
\end{solution}

\begin{solution}{exo:hf:on-chain-trading:6}
Its hedge is on a centralised exchange, outside the chain's transaction: one leg can execute without the other, leaving price risk and the risk of
inclusion failure.
\end{solution}

\begin{solution}{exo:hf:on-chain-trading:7}
Loss-versus-rebalancing rises to \$96\,300 over the week (the formula: \$112\,000, four times the base case); the hedged provider nets \$294\,000
instead of \$313\,000, with more fees from the larger arbitrage volume offsetting part of it.
\end{solution}

\begin{solution}{exo:hf:on-chain-trading:8}
\$11\,400 is the gross value before gas and the \omterm{def:hf:on-chain-trading:bid}{bundle bids}; after \$7\,100 of gas the searchers' net is \$4\,300, and in a competitive auction most of
that goes to builders: the searcher keeps its cost advantage over the next-best.
\end{solution}

\begin{solution}{pb:hf:on-chain-trading:1}
\begin{enumerate}
\item Maximise $g\Delta x\,y/(x+g\Delta x)-p\Delta x$: $g xy/(x+g\Delta x)^2=p$, so $\Delta x^\ast=(\sqrt{gxy/p}-x)/g$.
\item See \cref{def:hf:on-chain-trading:atomic}; CEX--DEX has an off-chain leg.
\item Arbitrage bots paying high fees and optimising latency to front-run users; priority gas auctions; consensus risk from ordering fees.
\item Searcher to builder with a bid; builder to proposer with a payment; the best block is proposed.
\item See \cref{def:hf:on-chain-trading:bid}.
\item $a+\frac{n-1}{n}(v-a)$; 93\%, 97\% and 99\%.
\item Its cost advantage over the second-best: \$6.7 of \$100 against one rival, \$1.0 against nineteen.
\item The same way: many searchers watch health factors; the best bid wins the liquidation.
\item See \cref{def:hf:on-chain-trading:hedged}; $\sigma^2V/8$ per unit time.
\item \$25\,800 to \$29\,500 a week against the formula's \$28\,000.
\item The price impact the noise traders paid, \$20\,000 to \$22\,500 a week; the noise traders.
\item It holds noise volume fixed, while real volume falls with the fee.
\item Builders receive 93\% of the winner's value with two searchers and 99\% with twenty; the hedged provider nets \$313\,000 on \$20 million in a
week at a 0.3\% fee, after \$26\,000 of loss-versus-rebalancing.
\item Lower costs than its rivals: better latency, inventory, gas.
\item The victim's loss (\$990) funds the attacker's gross (\$700), of which the builder takes most.
\item A mistrial in November 2025 in the Peraire-Bueno MEV-Boost case; the law's application unsettled.
\item Just-in-time liquidity.
\item Fees from realistic volume against $\sigma^2V/8$.
\item The value of its trade's impact would be auctioned back to it.
\item The builder, mostly.
\end{enumerate}
\end{solution}

\begin{solution}{iq:hf:on-chain-trading:1}
The pool is stale; the first arbitrage in the next block trades it back to within the fee of the new price; the searchers compete, and the builder
receives most of the value; the liquidity providers bear the loss.
\iqlookfor{arbitrage, auction, loss-versus-rebalancing.}
\end{solution}

\begin{solution}{iq:hf:on-chain-trading:2}
Arbitrage keeps the pool at the market price; for a constant-product pool the value is $2\sqrt{kP}$, whose second derivative in $P$ is
$-\frac12\sqrt k P^{-3/2}$; the loss rate is $-\frac12\sigma^2P^2V^{\prime\prime}=\frac{\sigma^2}{8}V$.
\iqlookfor{the pool value's concavity.}
\end{solution}

\begin{solution}{iq:hf:on-chain-trading:3}
Stream pending transactions and state; simulate candidate bundles against the latest state; estimate the value and competitors' bids; submit to
several builders before the deadline; track inclusion.
\iqlookfor{simulation against state, deadline, multiple builders.}
\end{solution}

\begin{solution}{iq:hf:on-chain-trading:4}
A naked position on the centralised exchange: unwind it at once, or resubmit the on-chain leg if still profitable; the system must treat inclusion
failure as a normal event.
\iqlookfor{\omterm{def:hf:index-and-basket-arbitrage:legging}{legging risk} on-chain.}
\end{solution}

\begin{solution}{iq:hf:on-chain-trading:5}
Bid below your value by less than you would against equals; if the competitor's value is surely higher, bid only when yours is close to it, or do
not bid.
\iqlookfor{asymmetric auctions.}
\end{solution}

\begin{solution}{iq:hf:on-chain-trading:6}
Against others bidding $\beta(v)=a+\frac{n-1}{n}(v-a)$, a bid $b$ wins with probability $\bigl(\frac{b-a}{\frac{n-1}{n}(b_{\max}-a)}\bigr)^{n-1}$
up to scale; maximising $(v-b)(b-a)^{n-1}$ gives $b=a+\frac{n-1}{n}(v-a)$.
\iqlookfor{the first-order condition.}
\end{solution}
